\section{Induksi pada Keterbagian dan Sifat Bilangan}

Induksi matematika juga digunakan secara luas untuk membuktikan sifat-sifat keterbagian (\textit{divisibility}) bilangan bulat. Pembuktian keterbagian melalui induksi sering menggunakan teknik pemfaktoran pada langkah induksi untuk mengisolasi faktor pembagi \cite{rosen2019}.

\subsection{Keterbagian oleh 3}

\begin{contoh}
\textbf{Teorema}: Untuk setiap bilangan bulat $n \geq 1$, $3 \mid (n^3 - n)$ (dibaca: 3 membagi $n^3 - n$).

\textbf{Bukti (Induksi):}

\textbf{Langkah Basis} ($n = 1$):
\[ 1^3 - 1 = 0 = 3 \cdot 0. \quad 3 \mid 0. \quad \checkmark \]

\textbf{Hipotesis Induksi}: Asumsikan untuk $n = k$:
\[ 3 \mid (k^3 - k), \text{ yaitu } k^3 - k = 3m \text{ untuk suatu bilangan bulat } m. \]

\textbf{Langkah Induksi}: Buktikan $3 \mid ((k+1)^3 - (k+1))$.
\begin{align*}
(k+1)^3 - (k+1) &= k^3 + 3k^2 + 3k + 1 - k - 1 \\
&= (k^3 - k) + 3k^2 + 3k \\
&= \underbrace{(k^3 - k)}_{= 3m \text{ (hip. ind.)}} + 3(k^2 + k) \\
&= 3m + 3(k^2 + k) \\
&= 3(m + k^2 + k)
\end{align*}

Karena $(m + k^2 + k)$ adalah bilangan bulat, maka $3 \mid ((k+1)^3 - (k+1))$. $\blacksquare$
\end{contoh}

\subsection{Keterbagian oleh 6}

\begin{contoh}
\textbf{Teorema}: Untuk setiap bilangan bulat $n \geq 1$, $6 \mid n(n+1)(n+2)$.

\textbf{Bukti (Induksi):}

\textbf{Langkah Basis} ($n = 1$):
\[ 1 \cdot 2 \cdot 3 = 6 = 6 \cdot 1. \quad 6 \mid 6. \quad \checkmark \]

\textbf{Hipotesis Induksi}: Asumsikan untuk $n = k$:
\[ 6 \mid k(k+1)(k+2), \text{ yaitu } k(k+1)(k+2) = 6m \text{ untuk suatu bilangan bulat } m. \]

\textbf{Langkah Induksi}: Buktikan $6 \mid (k+1)(k+2)(k+3)$.
\begin{align*}
(k+1)(k+2)(k+3) &= (k+1)(k+2) \cdot k + (k+1)(k+2) \cdot 3 \\
&= k(k+1)(k+2) + 3(k+1)(k+2) \\
&= \underbrace{k(k+1)(k+2)}_{= 6m \text{ (hip. ind.)}} + 3(k+1)(k+2)
\end{align*}

Suku pertama $= 6m$ yang habis dibagi 6. Untuk suku kedua, perhatikan bahwa $(k+1)(k+2)$ adalah hasil kali dua bilangan bulat berurutan, sehingga salah satunya pasti genap. Maka $(k+1)(k+2) = 2p$ untuk suatu bilangan bulat $p$, dan $3(k+1)(k+2) = 3 \cdot 2p = 6p$.

\[ (k+1)(k+2)(k+3) = 6m + 6p = 6(m + p) \]

Sehingga $6 \mid (k+1)(k+2)(k+3)$. $\blacksquare$
\end{contoh}

\subsection{Keterbagian Eksponensial}

\begin{contoh}
\textbf{Teorema}: Untuk setiap bilangan bulat $n \geq 0$, $3 \mid (4^n - 1)$.

\textbf{Bukti (Induksi):}

\textbf{Langkah Basis} ($n = 0$):
\[ 4^0 - 1 = 0 = 3 \cdot 0. \quad 3 \mid 0. \quad \checkmark \]

\textbf{Hipotesis Induksi}: Asumsikan untuk $n = k$:
\[ 4^k - 1 = 3m \text{ untuk suatu bilangan bulat } m, \text{ sehingga } 4^k = 3m + 1. \]

\textbf{Langkah Induksi}: Buktikan $3 \mid (4^{k+1} - 1)$.
\begin{align*}
4^{k+1} - 1 &= 4 \cdot 4^k - 1 \\
&= 4(3m + 1) - 1 & \text{(substitusi hip. ind.)} \\
&= 12m + 4 - 1 \\
&= 12m + 3 \\
&= 3(4m + 1)
\end{align*}

Karena $(4m + 1)$ bilangan bulat, maka $3 \mid (4^{k+1} - 1)$. $\blacksquare$
\end{contoh}
